\documentclass[10pt,a4paper]{amsart}
\usepackage{amsmath,amssymb,amsthm,graphicx}
\usepackage{verbatim}
\newenvironment{reference}{\comment}{\endcomment}
%\newenvironment{reference}{}{}
\newenvironment{slogan}{\comment}{\endcomment}
\newenvironment{history}{\comment}{\endcomment}

% For commutative diagrams we use Xy-pic
\usepackage[all]{xy}

% We use 2cell for 2-commutative diagrams.
\xyoption{2cell}
% Theorem environments.
%
\theoremstyle{plain}
\newtheorem{theorem}[subsection]{Theorem}
\newtheorem{proposition}[subsection]{Proposition}
\newtheorem{lemma}[subsection]{Lemma}

\theoremstyle{definition}
\newtheorem{definition}[subsection]{Definition}
\newtheorem{example}[subsection]{Example}
\newtheorem{exercise}[subsection]{Exercise}
\newtheorem{situation}[subsection]{Situation}

\theoremstyle{remark}
\newtheorem{remark}[subsection]{Remark}
\newtheorem{remarks}[subsection]{Remarks}

\numberwithin{equation}{subsection}

% Macros
%
\def\lim{\mathop{\mathrm{lim}}\nolimits}
\def\colim{\mathop{\mathrm{colim}}\nolimits}
\def\Spec{\mathop{\mathrm{Spec}}}
\def\Hom{\mathop{\mathrm{Hom}}\nolimits}
\def\Ext{\mathop{\mathrm{Ext}}\nolimits}
\def\SheafHom{\mathop{\mathcal{H}\!\mathit{om}}\nolimits}
\def\SheafExt{\mathop{\mathcal{E}\!\mathit{xt}}\nolimits}
\def\Sch{\mathit{Sch}}
\def\Mor{\mathop{\mathrm{Mor}}\nolimits}
\def\Ob{\mathop{\mathrm{Ob}}\nolimits}
\def\Sh{\mathop{\mathit{Sh}}\nolimits}
\def\NL{\mathop{N\!L}\nolimits}
\def\CH{\mathop{\mathrm{CH}}\nolimits}
\def\proetale{{pro\text{-}\acute{e}tale}}
\def\etale{{\acute{e}tale}}
\def\QCoh{\mathit{QCoh}}
\def\Ker{\mathop{\mathrm{Ker}}}
\def\Im{\mathop{\mathrm{Im}}}
\def\Coker{\mathop{\mathrm{Coker}}}
\def\Coim{\mathop{\mathrm{Coim}}}

% Boxtimes
%
\DeclareMathSymbol{\boxtimes}{\mathbin}{AMSa}{"02}

%
% Macros for moduli stacks/spaces
%
\def\QCohstack{\mathcal{QC}\!\mathit{oh}}
\def\Cohstack{\mathcal{C}\!\mathit{oh}}
\def\Spacesstack{\mathcal{S}\!\mathit{paces}}
\def\Quotfunctor{\mathrm{Quot}}
\def\Hilbfunctor{\mathrm{Hilb}}
\def\Curvesstack{\mathcal{C}\!\mathit{urves}}
\def\Polarizedstack{\mathcal{P}\!\mathit{olarized}}
\def\Complexesstack{\mathcal{C}\!\mathit{omplexes}}
% \Pic is the operator that assigns to X its picard group, usage \Pic(X)
% \Picardstack_{X/B} denotes the Picard stack of X over B
% \Picardfunctor_{X/B} denotes the Picard functor of X over B
\def\Pic{\mathop{\mathrm{Pic}}\nolimits}
\def\Picardstack{\mathcal{P}\!\mathit{ic}}
\def\Picardfunctor{\mathrm{Pic}}
\def\Deformationcategory{\mathcal{D}\!\mathit{ef}}
\usepackage{cleveref}
\crefname{originallemma}{original lemma}{original lemmas}
\Crefname{originallemma}{Original lemma}{Original lemmas}
\makeatletter
\newlabel{lemma-exterior-product-associative-dim-1}{{0FC8}{source}{}{sourcepointer-1}{}}
\newlabel{lemma-exterior-product-associative-dim-1@cref}{{[originallemma][0][]0FC8}{[1][0][]source}{}{}{}}
\makeatother
\crefname{originalreference}{original reference}{original references}
\Crefname{originalreference}{Original reference}{Original references}
\makeatletter
\newlabel{remark-restriction-bivariant}{{0F9Z}{source}{}{sourcepointer-2}{}}
\newlabel{remark-restriction-bivariant@cref}{{[originalreference][0][]0F9Z}{[1][0][]source}{}{}{}}
\makeatother
\crefname{originallemma}{original lemma}{original lemmas}
\Crefname{originallemma}{Original lemma}{Original lemmas}
\makeatletter
\newlabel{lemma-gysin-commutes}{{0FBP}{source}{}{sourcepointer-3}{}}
\newlabel{lemma-gysin-commutes@cref}{{[originallemma][0][]0FBP}{[1][0][]source}{}{}{}}
\makeatother
\crefname{originalreference}{original reference}{original references}
\Crefname{originalreference}{Original reference}{Original references}
\makeatletter
\newlabel{remark-commuting-exterior}{{0FBY}{source}{}{sourcepointer-4}{}}
\newlabel{remark-commuting-exterior@cref}{{[originalreference][0][]0FBY}{[1][0][]source}{}{}{}}
\makeatother
\title{The exterior product over a base of dimension one is associative}
\author{The Stacks Project Authors}
\date{stacks-project:chow.tex:local-snapshot-20261012; GNU FDL 1.2 or later}
\begin{document}
\maketitle
\noindent\textit{Editorial note.} All delivered mathematical text is the author's own, copied from the pinned Stacks Project chapter file chow.tex; no statement or proof step was rewritten, reordered or removed.\par\medskip
\noindent\textit{Editorial note.} Preamble provenance: the pinned chapter file opens with an input of the shared preamble and carries no preamble text of its own, so the project's single Stacks-wide preamble staged at sources/stacks-chow/preamble.tex is excerpted verbatim as two spans, lines 14-18 (Xy-pic setup) and lines 154-217 (the native theorem declarations and mathematical macros). Nothing was added, removed or redefined; the wrapper class and the amsmath/amssymb/amsthm/graphicx base packages are supplied by the assembler.\par\medskip
\noindent\textit{Editorial note.} The counted statement and proof contain no cross reference, no citation and no included file, so the delivered body is self contained.\par\medskip
\noindent\textit{Editorial note.} Canonical location of the counted proposition: Stacks tag 0FCA (sources/stacks-chow/pinned-original-source.tex line 14918).\par\medskip
\noindent\textit{Allowed standard prerequisites (editorial).}\begin{itemize}
\item Associativity of the exterior product over S in the dimension one setting. Same chapter, Lemma labelled lemma-exterior-product-associative-dim-1, cited at proof lines 14927-14929. Gives the associativity of the untwisted product that the proof transfers to the twisted one.
\item The bivariant classes of the diagonal immersions and their commutation. Same chapter, Lemma labelled lemma-gysin-commutes, cited at proof lines 14950-14957, with Remark labelled remark-restriction-bivariant for the restrictions. Produces the relation Delta\textasciicircum{}! o Delta\_12\textasciicircum{}! = Delta\textasciicircum{}! o Delta\_23\textasciicircum{}! used to reorder the product.
\item The product defined by the diagonal class. Same chapter, the construction preceding this lemma, invoked in the statement of the lemma. Fixes the meaning of alpha . beta as Delta\textasciicircum{}!(alpha x beta).
\item Commutation of the exterior product with the diagonal class. Same chapter, Remark labelled remark-commuting-exterior, cited at proof lines 14970-14975. Justifies the second and penultimate equalities of the chain.
\end{itemize}\medskip
\ifdefined\hypertarget\hypertarget{sourcepointer-1}{}\fi
\noindent\textit{Source reference.} Original lemma 0FC8: sources/stacks-chow/pinned-original-source.tex, source lines 14841--14841. Same-chapter pointer into the author environment lemma at source line 14841; invoked at proof line 14928 of the counted proof. It is retained only as a navigation pointer; every load-bearing step is inside the delivered span.\par\medskip
\ifdefined\hypertarget\hypertarget{sourcepointer-2}{}\fi
\noindent\textit{Source reference.} Original reference 0F9Z: sources/stacks-chow/pinned-original-source.tex, source lines 6184--6184. Same-chapter pointer into the author environment remark at source line 6184; invoked at proof line 14943 of the counted proof. It is retained only as a navigation pointer; every load-bearing step is inside the delivered span.\par\medskip
\ifdefined\hypertarget\hypertarget{sourcepointer-3}{}\fi
\noindent\textit{Source reference.} Original lemma 0FBP: sources/stacks-chow/pinned-original-source.tex, source lines 12206--12206. Same-chapter pointer into the author environment lemma at source line 12206; invoked at proof line 14949 of the counted proof. It is retained only as a navigation pointer; every load-bearing step is inside the delivered span.\par\medskip
\ifdefined\hypertarget\hypertarget{sourcepointer-4}{}\fi
\noindent\textit{Source reference.} Original reference 0FBY: sources/stacks-chow/pinned-original-source.tex, source lines 14261--14261. Same-chapter pointer into the author environment remark at source line 14261; invoked at proof line 14974 of the counted proof. It is retained only as a navigation pointer; every load-bearing step is inside the delivered span.\par\medskip
\begin{lemma}
\label{lemma-associative-dim-1}
The product defined above is associative. More precisely, with
$(S, \delta)$ as above, let $X$ be smooth over $S$,
let $Y, Z, W$ be schemes locally of finite type over $X$, let
$\alpha \in \CH_*(Y)$, $\beta \in \CH_*(Z)$, $\gamma \in \CH_*(W)$.
Then $(\alpha \cdot \beta) \cdot \gamma =
\alpha \cdot (\beta \cdot \gamma)$ in $\CH_*(Y \times_X Z \times_X W)$.
\end{lemma}

\begin{proof}
By Lemma \ref{lemma-exterior-product-associative-dim-1} we have
$(\alpha \times \beta) \times \gamma =
\alpha \times (\beta \times \gamma)$ in $\CH_*(Y \times_S Z \times_S W)$.
Consider the closed immersions
$$
\Delta_{12} : X \times_S X \longrightarrow X \times_S X \times_S X,
\quad (x, x') \mapsto (x, x, x')
$$
and
$$
\Delta_{23} : X \times_S X \longrightarrow X \times_S X \times_S X,
\quad (x, x') \mapsto (x, x', x')
$$
Denote $\Delta_{12}^!$ and $\Delta_{23}^!$ the corresponding bivariant
classes; observe that $\Delta_{12}^!$ is the restriction
(Remark \ref{remark-restriction-bivariant}) of $\Delta^!$
to $X \times_S X \times_S X$ by the map $\text{pr}_{12}$ and that
$\Delta_{23}^!$ is the restriction of $\Delta^!$
to $X \times_S X \times_S X$ by the map $\text{pr}_{23}$.
Thus clearly the restriction of $\Delta_{12}^!$ by $\Delta_{23}$
is $\Delta^!$ and the restriction of $\Delta_{23}^!$ by $\Delta_{12}$ is
$\Delta^!$ too. Thus by Lemma \ref{lemma-gysin-commutes} we have
$$
\Delta^! \circ \Delta_{12}^! =
\Delta^! \circ \Delta_{23}^! 
$$
Now we can prove the lemma by the following sequence of equalities:
\begin{align*}
(\alpha \cdot \beta) \cdot \gamma
& =
\Delta^!(\Delta^!(\alpha \times \beta) \times \gamma) \\
& =
\Delta^!(\Delta_{12}^!((\alpha \times \beta) \times \gamma)) \\
& =
\Delta^!(\Delta_{23}^!((\alpha \times \beta) \times \gamma)) \\
& =
\Delta^!(\Delta_{23}^!(\alpha \times (\beta \times \gamma)) \\
& =
\Delta^!(\alpha \times \Delta^!(\beta \times \gamma)) \\
& =
\alpha \cdot (\beta \cdot \gamma)
\end{align*}
All equalities are clear from the above except perhaps
for the second and penultimate one. The equation
$\Delta_{23}^!(\alpha \times (\beta \times \gamma)) =
\alpha \times \Delta^!(\beta \times \gamma)$ holds by
Remark \ref{remark-commuting-exterior}. Similarly for the second
equation.
\end{proof}

\end{document}
